% =====================================================================
%  MATH300 Project Description -- Nonlinear ODEs
% =====================================================================
\documentclass[11pt,a4paper]{article}

\usepackage[T1]{fontenc}
\usepackage[utf8]{inputenc}
\usepackage{lmodern}
\usepackage[margin=2.5cm]{geometry}
\usepackage{enumitem}
\usepackage{booktabs}
\usepackage{tabularx}
\usepackage{array}
\usepackage{amsmath,amssymb}
\usepackage{xcolor}
\usepackage[colorlinks=true,urlcolor=blue,linkcolor=black,citecolor=black]{hyperref}

\setlist{itemsep=2pt, topsep=4pt}
\setlength{\parindent}{0pt}
\setlength{\parskip}{6pt}

% Placeholder macro: blue text in square brackets.
% To hide the colour in the final version, change {blue} to {black}.
\newcommand{\ph}[1]{\textcolor{blue}{[#1]}}

% Marker for optional bonus tasks.
\newcommand{\bonus}{\textbf{(Bonus)}}

\newcolumntype{L}{>{\raggedright\arraybackslash}X}

\begin{document}

% ---------------------------------------------------------------------
%  Title (no logo, no instructor name)
% ---------------------------------------------------------------------
\begin{center}
  {\large\bfseries MATH300: Mathematical Computing \textperiodcentered{} Fall 2026}\\[6pt]
  {\Large\bfseries Project: Nonlinear Dynamics in Continuous Time}
\end{center}
\vspace{6pt}

Groups must be formed, and project titles submitted on Canvas, by
\textbf{October 26, 2026}. The project
has two graded deliverables and an oral exam: a 3--5 page report for
Project~2 (due \textbf{November 20, 2026}), a poster for the final project
(due \textbf{December 11, 2026}), and an oral exam held between November 30
and December 17, 2026, 3:30\,pm. Read the general rules below before you
start; they apply throughout the project.

% ---------------------------------------------------------------------
\section*{General remarks}
% ---------------------------------------------------------------------
\begin{itemize}
  \item Projects are set for a group of \textbf{two or three} students.
  \item Form your group by \textbf{October 26, 2026}, and submit your
        project title as a group assignment on \textbf{Canvas} by the same
        date.
  \item Every group member must carry an \textbf{equal portion} of the coding
        and of the report and poster preparation.
  \item The scope can be adjusted to the number of participants, their
        interests and prior knowledge.
  \item Only \textbf{Python} may be used for programming.
  \item Implementation, documentation, validation and demonstration examples
        are all part of the project.
  \item Within these rules you are free to choose the implementation, code
        structure and exact content. You can discuss this with me at any time.
\end{itemize}

% ---------------------------------------------------------------------
\section*{Version control and GitHub}
% ---------------------------------------------------------------------
Each group keeps \textbf{one GitHub repository for the whole project}
(Project~2 and the final project), and every member commits from their own
GitHub account.
\begin{itemize}
  \item Use git from the very beginning. Commits should be small, consistent
        and clearly described. Avoid large commits.
  \item Each group member uploads their own code using \textbf{their own
        GitHub account}. Commits from a shared or borrowed account will not
        be credited to the other members.
  \item The repository name is up to your group.
  \item If the repository is private, add my GitHub account
        \textbf{\texttt{say-yas}} as a collaborator.
  \item The repository must contain a README describing the project, how to
        install and run the code, and the AI-use statement (see below).
\end{itemize}

% ---------------------------------------------------------------------
\section*{Use of AI tools}
% ---------------------------------------------------------------------
\textbf{AI-generated code is not accepted in any form.} Code produced by a
chat assistant or code generator (e.g.\ ChatGPT, Claude, Gemini) and copied,
pasted or adapted into your project will be treated as work that is not your
own.
\begin{itemize}
  \item \textbf{Permitted:} inline coding assistance such as GitHub Copilot
        autocompletion while you write the code yourself, \textbf{provided its
        use is acknowledged}.
  \item \textbf{Acknowledgement:} include an AI-use statement in the README
        and in the report that names the tool(s) used and the parts of the
        code (files or functions) where they assisted. If no AI tools were
        used, state this explicitly.
  \item \textbf{Reports and posters must not be generated by AI} in any
        form: their text, figures and layout must be your own work. If AI
        tools were used for any other non-code purpose (e.g.\ checking
        grammar or searching the literature), describe this in the same
        statement.
  \item You are responsible for every line in your repository. In the oral
        exam you must be able to explain and modify any part of the code you
        submitted.
\end{itemize}

% ---------------------------------------------------------------------
\section*{Deliverables and deadlines}
% ---------------------------------------------------------------------
Submit each deliverable following the details provided on the
\textbf{Canvas} page. The written report and poster are handed in on
\textbf{Crowdmark}, and the link to your GitHub repository is submitted
separately on \textbf{Canvas}.

\begin{center}
\small
\renewcommand{\arraystretch}{1.35}
\begin{tabularx}{\textwidth}{|>{\raggedright\arraybackslash}p{2.5cm}|>{\raggedright\arraybackslash}p{3.0cm}|>{\raggedright\arraybackslash}p{3.0cm}|L|}
\hline
\textbf{Deliverable} & \textbf{Due} & \textbf{Submission} & \textbf{Notes} \\
\hline\hline
Group and project title & Monday, October 26, 2026 &
  Group assignment on Canvas &
  Group of two or three members; submit the project title. \\
\hline
Project~2 report & Friday, November 20, 2026 &
  Report (PDF) on Crowdmark; repository link on Canvas &
  3--5 pages; the page limit \textbf{includes} references, figures and
  appendices. \\
\hline
Final project poster & Friday, December 11, 2026 &
  Poster (PDF) on Crowdmark; repository link on Canvas &
  36 inches high by 48 inches wide (landscape). Target audience:
  undergraduate classmates with a basic background in math, physics or
  engineering. \\
\hline
Oral exam & Between November 30 and December 17, 2026, 3:30\,pm &
  Individual meeting; details discussed in class &
  30 minutes per student. Your poster must already be submitted on
  Crowdmark before your exam. The last possible slot ends December 17 at
  3:30\,pm. \\
\hline
\end{tabularx}
\end{center}

\subsection*{Oral exam}
The oral exam is individual and lasts 30 minutes per student. It takes place
between November 30 and December 17, 2026, 3:30\,pm. The oral exams and how
to arrange them will be discussed in class before the Thanksgiving break.

Before your oral exam, your poster must also be submitted on
\textbf{Crowdmark}; that is, you must have finished both Project~2 and the
final project. In the meeting you will:
\begin{itemize}
  \item Present your results and algorithms.
  \item Possibly live-code and run parts of your algorithms to generate
        outputs on the spot.
  \item Answer questions about the theory, your implementation and your
        results. Each member should be able to explain the whole project, not
        only their own parts.
\end{itemize}
Bring a laptop with the repository cloned and a working environment, so that
the code runs without additional setup.

% ---------------------------------------------------------------------
\section*{Grading}
% ---------------------------------------------------------------------
\begin{center}
\small
\renewcommand{\arraystretch}{1.35}
\begin{tabularx}{0.8\textwidth}{|L|>{\raggedleft\arraybackslash}p{3.2cm}|}
\hline
\textbf{Component} & \textbf{Points} \\
\hline\hline
Group and project title (on time) & 10 \\
\hline
Project~2 report & 100 \\
\hline
Final project poster & 100 \\
\hline
Oral exam (individual) & 50 \\
\hline
Bonus: extra steps beyond the required tasks & up to 10 \\
\hline
\end{tabularx}
\end{center}

The Project~2 report and the poster are each graded out of 100. What matters
most is:
\begin{itemize}
  \item \textbf{Coherence:} a clear story from the question, through the
        methods, to the results and conclusions.
  \item \textbf{Correctness:} of the mathematics, the code and the results.
  \item \textbf{Depth of analysis and discussion:} explain what your results
        mean, why they look the way they do, and what their limitations are,
        rather than only showing them.
  \item \textbf{Professional presentation:} figures with labeled axes and
        readable font sizes, captions and legends, a clean and readable
        layout, and proper references.
\end{itemize}

A complete, correct and well-discussed
solution of the required tasks can earn full marks. Tasks marked
\bonus{} and the suggestions under ``Going further'' are optional; exploring
such extra steps can earn \textbf{up to 10 bonus points} in total.

% ---------------------------------------------------------------------
\section*{Requirements for the report and the poster}
% ---------------------------------------------------------------------
Both the report and the poster must contain a clear description of the
project, references, properly labeled figures and an author-contributions
statement.

\subsection*{Content}
\begin{itemize}
  \item A clear and concise description of the problem, the methods and the
        results.
  \item The AI-use statement described above.
  \item The report should not contain any significant amount of code; refer
        to the repository instead.
\end{itemize}

\subsection*{Poster}
\begin{itemize}
  \item The poster must be \textbf{self-contained}: a reader should be able to
        follow it without the report, the code or a spoken explanation.
  \item You may be selective in what you show and how you narrate your story,
        so that an undergraduate audience with a basic background in math,
        physics or engineering can follow it.
  \item You may choose your own title for the poster; be creative.
\end{itemize}

\subsection*{References}
\begin{itemize}
  \item Include a list of references at the end of the report and the poster.
  \item Cite sources clearly where they are used in the text (papers, books,
        software, code repositories and data sets).
  \item Use one consistent citation style throughout.
\end{itemize}

\subsection*{Figures}
\begin{itemize}
  \item Every plot must have labeled $x$ and $y$ axes (and $z$ axis for 3D
        plots or color bars), including units where applicable.
  \item Axis labels, tick labels and legends must have a font size similar to
        the surrounding text, so they remain readable without zooming.
  \item Each figure needs a legend when it shows more than one data series,
        and a caption that states what is shown.
  \item Use log scales where appropriate (e.g.\ for convergence plots) and say
        so in the caption.
\end{itemize}

\subsection*{Author contributions}
Each group must include an author-contributions statement in the report and
the README that states:
\begin{itemize}
  \item Who did what: which parts of the code, experiments, analysis, writing
        and poster each member worked on.
  \item To what extent: the approximate share of each member in each part
        (e.g.\ 50\,\% / 50\,\%).
  \item Who took the lead in which parts.
\end{itemize}
The coding and the report and poster preparation must be shared equally
between the members. The commit history should be consistent with the stated
contributions.

% ---------------------------------------------------------------------
\section*{Codebase requirements}
% ---------------------------------------------------------------------
\begin{itemize}
  \item Well structured and documented, both in the code (docstrings,
        comments) and in the repository (README).
  \item Validated: show that your implementation produces correct results,
        e.g.\ on examples with known solutions.\end{itemize}

% ---------------------------------------------------------------------
\section*{Project description}
% ---------------------------------------------------------------------

\subsection*{Goal}
Many systems in physics, chemistry, biology and engineering change
continuously in time and are described by nonlinear ordinary differential
equations (ODEs): a swinging pendulum, a predator--prey population, a chemical
oscillator, or the Lorenz model of convection. In this project you will choose
\textbf{at least 3 systems} from the list in \cite{list}, solve them
numerically with SciPy's ODE solver \texttt{solve\_ivp} \cite{scipy}, and
explore their behavior with phase portraits.

\subsection*{Background}
A system of first-order ODEs has the form
\begin{equation*}
  \frac{d\mathbf{x}}{dt} = \mathbf{f}(\mathbf{x}, t), \qquad
  \mathbf{x}(0) = \mathbf{x}_0,
\end{equation*}
where $\mathbf{x}(t)$ is the state of the system at time $t$. A second-order
equation can always be rewritten in this form. For example, the Van der Pol
oscillator $x'' - \mu(1 - x^2)x' + x = 0$ becomes the two-dimensional system
\begin{equation*}
  x' = y, \qquad y' = \mu(1 - x^2)\,y - x .
\end{equation*}
A famous three-dimensional example is the Lorenz system \cite{lorenz},
\begin{equation*}
  x' = \sigma(y - x), \qquad y' = x(\rho - z) - y, \qquad z' = xy - \beta z,
\end{equation*}
whose solutions, for $\sigma = 10$, $\rho = 28$ and $\beta = 8/3$, never settle
down and never repeat: they are chaotic.

Most nonlinear ODEs cannot be solved by hand, so the solution is approximated
numerically at a sequence of times. In this project you use SciPy's
\texttt{solve\_ivp} \cite{scipy}. By default it uses a Runge--Kutta method
that chooses its step size automatically, so that the estimated error in each
step stays below given tolerances (the arguments \texttt{rtol} and
\texttt{atol}). Smaller tolerances give more accurate solutions but take
longer to compute. You still have to choose these settings, and the time
interval and initial states, sensibly.

The \emph{phase space} of a system is the set of all its possible states, for
example the plane of all points $(x, y)$ for the Van der Pol oscillator. A
\emph{phase portrait} draws solutions in phase space instead of against time:
for a two-dimensional system you plot $y(t)$ against $x(t)$, so time is not
shown on any axis. Each solution appears as a continuous curve, called a
trajectory, and drawing trajectories from many initial states shows
where they go in the long run: to an equilibrium point, around a closed loop
(a limit cycle, as for the Van der Pol oscillator), or onto a strange
attractor (as for the Lorenz system). A good introduction to all of these
ideas is \cite{strogatz}.

\subsection*{Tasks for Project 2 (report, due November 20, 2026)}
\begin{enumerate}
  \item Choose \textbf{at least 2 systems} from \cite{list}. Include at least
        one two-dimensional system (such as the Van der Pol oscillator, the
        pendulum, the Lotka--Volterra equations or the FitzHugh--Nagumo model)
        and at least one three-dimensional system (such as the Lorenz or the
        R\"ossler system). Write each one as a system of first-order ODEs,
        describe what it models and the meaning of its parameters, and cite
        your sources.
  \item Write the right-hand side of each system as a Python function and
        solve it with \texttt{solve\_ivp} \cite{scipy}. Plot the solutions
        against time for a few initial states and parameter values.
  \item Check the accuracy of the solver on an equation with a known exact
        solution, for example the logistic equation $f' = f(1 - f)$ from
        \cite{list}. Plot the error at a fixed final time against the
        tolerance \texttt{rtol} on a log--log scale, and discuss how the
        tolerance affects the accuracy.
  \item \bonus\ Compare different methods offered by \texttt{solve\_ivp}
        (for example \texttt{RK45}, \texttt{DOP853} and \texttt{Radau}) in
        accuracy and computing time.
\end{enumerate}

\subsection*{Tasks for the final project (poster, due December 11, 2026)}
\begin{enumerate}
  \item Draw phase portraits of your systems (for the two-dimensional
        systems, trajectories from many initial states in the plane; for the
        three-dimensional systems, a 3D plot or two-dimensional projections),
        and compare them for several values of one parameter of each system
        (for example $\mu$ for the Van der Pol oscillator or $\rho$ for the
        Lorenz system).
  \item Discuss what you observe: what is the long-term behavior of each
        system, and how does it change with the parameter? What do your
        systems have in common, and how do they differ? What are the limitations of solving nonlinear
        ODEs numerically (accuracy, long simulation times, chaotic
        systems)?
  \item \bonus\ Show sensitivity to initial conditions for your
        three-dimensional system: start two solutions very close together
        and plot the distance between them against $t$ (log scale).
\end{enumerate}

\subsection*{Going further}
These optional steps count toward the bonus points. For example, you can find
the equilibrium points of your systems and check their stability, or compute
a Poincar\'e section of a three-dimensional system.

% ---------------------------------------------------------------------
\begin{thebibliography}{9}
  \bibitem{list} Wikipedia. \emph{List of nonlinear ordinary differential
    equations}.
    \url{https://en.wikipedia.org/wiki/List_of_nonlinear_ordinary_differential_equations}
  \bibitem{lorenz} E.~N.~Lorenz (1963). Deterministic nonperiodic flow.
    \emph{Journal of the Atmospheric Sciences}, 20, 130--141.
    \url{https://doi.org/10.1175/1520-0469(1963)020<0130:DNF>2.0.CO;2}
  \bibitem{strogatz} S.~H.~Strogatz (2015). \emph{Nonlinear Dynamics and
    Chaos}, 2nd edition. Westview Press.
  \bibitem{scipy} SciPy. \emph{scipy.integrate.solve\_ivp}.
    \url{https://docs.scipy.org/doc/scipy/reference/generated/scipy.integrate.solve_ivp.html}
\end{thebibliography}

\end{document}
